\section{Independent Set}
\label{sec:mwis}

In this section, we show the bound of \Cref{thm:main_is}. The class of graphs we use are inspired by the classic reduction from CNF-sat to Maximum Weight Independent Set.
Fix $k > 0$ and let $q = 4\binom{k}{2}$. Consider all possible 2-CNF clauses on $k$ variables: there are exactly $q$ of those. We can think of each clause as a tuple $(a, b, n^a, n^b)$ where $a, b \in [k]$ represents an unordered pair of variables and $n^a, n^b \in \{0, 1\}$ represents whether $a$-th and $b$-th variable is negated in the clause. We define the undirected graph $G_k$ in the following way. We put
\begin{equation*}
    V(G_k) = \{v_{i,j} \mid i \in [k], j \in \{0, \dots, 2q - 1\}\} \cup
             \{ w_j \mid j \in [q] \}.
\end{equation*}
The vertices $v_{i,j}$ represent the literals of the CNF formula: $i$ denotes the variable and the parity of $j$ determines whether it is negated or not.
Every literal has $q$ copies, one for each clause. The vertices $w_j$ represent all possible 2-CNF clauses. For the sake of convenience, we will use $w_{a,b,n^a,n^b}$ to denote the vertex corresponding to the clause represented by $a,b,n^a,n^b$ as described above. We put $V_{i,r} = \{ v_{i, 2j + r} \mid j \in \{0, \dots q - 1\} \}$ for $i \in [k], r \in \{0, 1\}$, that is, even and odd vertices representing the $i$-th variable. We also put $V_i = V_{i,0} \cup V_{i,1}$.

We put the edges accordingly:
\begin{equation*}
    E(G_k) = \{ (v_{i,j-1}, v_{i,j}) \mid i \in [k], j \in [2q - 1]\} \cup
             \{ (w_i, l_{i,j} \mid i \in [q], j \in \{1, 2\} \}
\end{equation*}
where $l_{i,1} = v_{a_i, 2(i - 1) + n^a_i}$ and $l_{i,2} = v_{b_i, 2(i - 1) + n^b_i}$ where $w_i = w_{a_i, b_i, n^a_i, n^b_i}$. That is, $l_{i,j}$ correspond to the literals of the clause represented by $w_i$. Because there is a copy $v_{i,j}$ of each literal for every clause, every $v_{i,j}$ is connected at most to one vertex $w_{1 + \lfloor j/2 \rfloor}$. The degree of each $w_i$ is exactly $2$.

\begin{figure}[h]
    \centering
    \label{fig:mwis}
    \begin{tikzpicture}
		\node [style=node] (0) at (-6, 1) [label={$v_{1,0}$}] {};
		\node [style=node-is] (1) at (-5, 1) [label={$v_{1,1}$}] {};
		\node [style=node] (2) at (-4, 1) {};
		\node [style=node-is] (3) at (-3, 1) {};
		\node [style=node] (4) at (-2, 1) {};
		\node [style=node-is] (5) at (-1, 1) {};
		\node [style=node] (6) at (0, 1) {};
		\node [style=node-is] (7) at (1, 1) {};
		\node [style=node] (8) at (2, 1) {};
		\node [style=node-is] (9) at (3, 1) {};
		\node [style=node-is] (10) at (-6, 0) [label={$v_{2,0}$}] {};
		\node [style=node] (11) at (-5, 0) [label={$v_{2,1}$}] {};
		\node [style=node-is] (12) at (-4, 0) {};
		\node [style=node] (13) at (-3, 0) {};
		\node [style=node-is] (14) at (-2, 0) {};
		\node [style=node] (15) at (-1, 0) {};
		\node [style=node-is] (16) at (0, 0) {};
		\node [style=node] (17) at (1, 0) {};
		\node [style=node-is] (18) at (2, 0) {};
		\node [style=node] (19) at (3, 0) {};
		\node [style=node] (20) at (-6, -1) [label={$v_{3,0}$}] {};
		\node [style=node-is] (21) at (-5, -1) [label={$v_{3,1}$}] {};
		\node [style=node] (22) at (-4, -1) {};
		\node [style=node-is] (23) at (-3, -1) {};
		\node [style=node] (24) at (-2, -1) {};
		\node [style=node-is] (25) at (-1, -1) {};
		\node [style=node] (26) at (0, -1) {};
		\node [style=node-is] (27) at (1, -1) {};
		\node [style=node] (28) at (2, -1) {};
		\node [style=node-is] (29) at (3, -1) {};
		\node [style=node] (30) at (-5.5, 2.25) [label={$w_{1,2,0,0}$}] {};
		\node [style=node-is] (31) at (-3.5, 2.25) [label={$w_{1,2,0,1}$}] {};
		\node [style=node] (32) at (-1.5, 2.25) [label={$w_{1,2,1,0}$}] {};
		\node [style=node] (33) at (0.5, 2.25) [label={$w_{1,2,1,1}$}] {};
		\node [style=node-is] (34) at (2.5, 2.25) [label={$w_{1,3,0,0}$}] {};
		\node [style=none] (35) at (4, 1) {};
		\node [style=none] (36) at (4, 0) {};
		\node [style=none] (37) at (4, -1) {};
		\node [style=none] (38) at (4.5, 0) {$\dots$};
		\node [style=none] (39) at (5, 1) {};
		\node [style=none] (40) at (5, 0) {};
		\node [style=none] (41) at (5, -1) {};
		\node [style=node] (42) at (6, 1) {};
		\node [style=node-is] (43) at (7, 1) {};
		\node [style=node] (44) at (7, 0) {};
		\node [style=node-is] (45) at (6, 0) {};
		\node [style=node] (46) at (6, -1) {};
		\node [style=node-is] (47) at (7, -1) {};
		\node [style=node-is] (48) at (6.5, 2.25) [label={$w_{2,3,1,0}$}] {};
		\node [style=node] (49) at (8, 1) [label={$v_{1,22}$}] {};
		\node [style=node-is] (50) at (9, 1) [label={$v_{1,23}$}] {};
		\node [style=node] (51) at (9, 0) [label={$v_{2,23}$}] {};
		\node [style=node-is] (52) at (8, 0) [label={$v_{2,22}$}] {};
		\node [style=node] (53) at (8, -1) [label={$v_{3,22}$}] {};
		\node [style=node-is] (54) at (9, -1) [label={$v_{3,23}$}] {};
		\node [style=node] (55) at (8.5, 2.25) [label={$w_{2,3,1,1}$}] {};
        
		\draw [style=undirected] (0) to (1);
		\draw [style=undirected] (1) to (2);
		\draw [style=undirected] (2) to (3);
		\draw [style=undirected] (3) to (4);
		\draw [style=undirected] (4) to (5);
		\draw [style=undirected] (5) to (6);
		\draw [style=undirected] (6) to (7);
		\draw [style=undirected] (7) to (8);
		\draw [style=undirected] (8) to (9);
		\draw [style=undirected] (10) to (11);
		\draw [style=undirected] (11) to (12);
		\draw [style=undirected] (12) to (13);
		\draw [style=undirected] (13) to (14);
		\draw [style=undirected] (14) to (15);
		\draw [style=undirected] (15) to (16);
		\draw [style=undirected] (16) to (17);
		\draw [style=undirected] (17) to (18);
		\draw [style=undirected] (18) to (19);
		\draw [style=undirected] (20) to (21);
		\draw [style=undirected] (21) to (22);
		\draw [style=undirected] (22) to (23);
		\draw [style=undirected] (23) to (24);
		\draw [style=undirected] (24) to (25);
		\draw [style=undirected] (25) to (26);
		\draw [style=undirected] (26) to (27);
		\draw [style=undirected] (27) to (28);
		\draw [style=undirected] (28) to (29);
		\draw [style=undirected] (0) to (30);
		\draw [style=undirected] (10) to (30);
		\draw [style=undirected] (31) to (2);
		\draw [style=undirected] (31) to (13);
		\draw [style=undirected] (32) to (14);
		\draw [style=undirected] (32) to (5);
		\draw [style=undirected] (33) to (7);
		\draw [style=undirected] (33) to (17);
		\draw [style=undirected] (34) to (8);
		\draw [style=undirected] (34) to (28);
		\draw [style=undirected] (9) to (35.center);
		\draw [style=undirected] (19) to (36.center);
		\draw [style=undirected] (29) to (37.center);
		\draw [style=undirected] (39.center) to (42);
		\draw [style=undirected] (40.center) to (45);
		\draw [style=undirected] (41.center) to (46);
		\draw [style=undirected] (42) to (43);
		\draw [style=undirected] (45) to (44);
		\draw [style=undirected] (46) to (47);
		\draw [style=undirected] (55) to (51);
		\draw [style=undirected] (55) to (54);
		\draw [style=undirected] (49) to (50);
		\draw [style=undirected] (52) to (51);
		\draw [style=undirected] (53) to (54);
		\draw [style=undirected] (48) to (44);
		\draw [style=undirected] (48) to (46);
		\draw [style=undirected] (43) to (49);
		\draw [style=undirected] (44) to (52);
		\draw [style=undirected] (47) to (53);
\end{tikzpicture}

    \caption{The graph $G_3$. Red vertices belong to the canonical solution for valuation $\rho(1) = 1, \rho(2) = 0, \rho(3) = 1$.}
\end{figure}

It is easy to see that the graph $G_k$ has pathwidth at most $k + 1$: We can start with bags $\{v_{i,j}:i \in [k]\}$ for $j \in \{0,\ldots,2q-1\}$ and add the sets $N[w_j]$ to them in the natural way. Proving the following lemma will immediately show \Cref{thm:main_is}.

\begin{lemma}\label{lem:mwis}
	For any $k \geq 1$, any tropical circuit calculating $IS_{G_k}$ uses at least $2^k / 3$ gates.
\end{lemma}

The rest of this section is dedicated to proving this lemma. Fix $k > 0$ and put $G := G_k$.
The proof will be done in two steps. First, we will show a combinatorial structure of so-called rectangles in our described graph. Then, we will relate this notion to the polynomials calculated by tropical circuits and show the desired lower bound as a consequence.

The \emph{canonical solution $I_\rho$ given by the assignment $\rho : [k] \to \{0, 1\}$} is defined as
\begin{equation*}
	I_\rho = \{ v_{i, 2j + \rho(i)} \mid i \in [k], j \in \{0, \dots, q-1\} \} \cup
	         \{ w_i \mid \rho(a_i) \neq n^a_i \ \text{and}\ \rho(b_i) \neq n^b_i \}.
\end{equation*}
That is, we pick vertices $v_{i,j}$ for all $j$ even or all $j$ odd depending on $\rho(i)$ and include all vertices $w_i$ corresponding to clauses not satisfied by assignment $\rho$. Clearly, there are $2^k$ canonical solutions, one per each assignment $\rho$, and every canonical solution is an independent set of $G$.

We will say that a pair $\cA, \cB \subseteq 2^{V(G)}$ of families of independent sets of $G$ forms a \emph{rectangle} in $G$ (denoted $\cA \cdot \cB$) if and only if every pair of sets $A \in \cA, B \in \cB$ are disjoint and non-adjacent (in particular, $A \cup B$ forms an independent set in $G$). We refer to the families $\cA, \cB$ as \emph{sides} of the rectangle. We will say that a rectangle $\cA \cdot \cB$ \emph{contains} an independent set $I$ if $I = A \cup B$ for some $A \in \cA, B \in \cB$. We will say that a set $A \in \cA$ is \emph{useful} if there exists $B \in \cB$ such that $A \cup B$ is a canonical solution. We define a set $B \in \cB$ being useful in an analogous way.

For simplicity, the notion of a rectangle introduced here is defined in terms of families of independent sets of $G$. This concept, however, will be crucial in showing all of the subsequent bounds, hence, in later sections we will redefine rectangles in terms of tropical polynomials, as in the statement of \Cref{lem:tc_decomposition}, in order for the definition to be more general.


\subsection{Thin rectangles}

The key combinatorial property of the circuits calculating $IS_G$ is expressed via the following lemma.
\begin{lemma}\label{lem:is_thin_rects}
	Let $\cA \cdot \cB$ be a rectangle in $G$ that contains at least one canonical solution $I_\rho$. Then, either $\cA$ or $\cB$ contain at most one useful set.
\end{lemma}

For the rest of this subsection, we fix a rectangle $\cA \cdot \cB$, and focus on proving the lemma via the following series of claims.
\begin{claim}
    Let $I_1, I_2$ be two different useful sets belonging to the same side of $\cA \cdot \cB$. Then, there exists $i \in [k], j \in \{0, \dots, 2q - 1\}$ such that $v_{i, j}$ belongs to exactly one of $I_1, I_2$.
\end{claim}
\begin{proof}
    W.l.o.g. we assume that $I_1, I_2 \in \cA$. If both sets have different intersection with some $V_i$, then the claim trivially follows, hence w.l.o.g. we assume that for some $l \in [q]$, we have $w_l \in I_1$ and $w_l \not\in I_2$.
    
    Since $I_2$ is useful, we have $J_2 \in \cB$ such that $I_2 \cup J_2$ is a canonical solution. Since $I_1$ and $J_2$ are disjoint, we have $w_l \not\in I_2 \cup J_2$, hence $v_{i,j} \in I_2 \cup J_2$ for some neighbor $v_{i, j}$ of $w_l$. As $I_1$ and $J_2$ are non-adjacent, we have $v_{i,j} \not\in J_2$, hence $v_{i,j} \in I_2$. As $I_1$ is an independent set of $G$, we have $v_{i,j} \not\in I_1$, which finishes the proof of the claim.
\end{proof}

\begin{claim}
    Let $I_1, I_2$ be two different useful sets belonging to the same side of $\cA \cdot \cB$. Then, there exists $i \in [k], r \in \{0,1\}$ such that $V_{i,r} \subseteq I_1$ and $V_{i, 1-r} \subseteq I_2$.
\end{claim}
\begin{proof}
    Assume w.l.o.g. that $I_1, I_2 \in \cA$.
    By the previous claim, we have some $i \in [k], j \in \{0, \dots, 2q - 1\}$ such that w.l.o.g. $v_{i,j} \in I_1$ and $v_{i,j} \not\in I_2$. Put $r = j \mod 2$. Since $I_1, I_2$ are useful, we have the corresponding $J_1, J_2 \in \cB$. By contradiction, let $j' \in \{0, \dots, 2q - 1\}$ be an index minimizing $|j - j'|$ for which $v_{i,j'}$ does not belong to the expected set, i.e., such that either $j' \mod 2 = r$ and $v_{i,j'} \not\in I_1$ or $j' \mod 2 \neq r$ and $v_{i,j'} \not\in I_2$.

    Consider the case $j' \mod 2 = r$, the other one will be symmetric. As $v_{i,j} \in I_1 \cup J_1$, we also have $v_{i,j'} \in I_1 \cup J_1$. By minimality, we have that either $v_{i,j'+1}$ or $v_{i,j'-1}$ is in $I_2$. If $v_{i,j'} \not\in I_1$ then $v_{i,j'} \in J_1$, but this would imply that $I_2 \cup J_1$ is not an independent set, which is a contradiction.
\end{proof}

\begin{claim}
    Let $I_1, I_2 \in \cA$ be two different useful sets and let $J_1 \in \cB$ be such that $I_1 \cup J_1$ is a canonical solution. Then $J_1$ is the only useful set in $\cB$.
\end{claim}
\begin{proof}
    Assume by contradiction that we have useful $J_2 \in \cB$ different to $J_1$. Applying the previous claim to both sides, we obtain $i_A, i_B \in [k]$ and $r_A, r_B \in \{0,1\}$ such that: $V_{i_A,r_A} \subseteq I_1, V_{i_A, 1 - r_A} \subseteq I_2, V_{i_B, r_B} \subseteq J_1$ and $V_{i_B, 1 - r_B} \subseteq J_2$. Naturally, we have $i_A \neq i_B$.

    Pick $w = w_{i_A, i_B, 1 - r_A, 1 - r_B}$. That is,  for some $j \in \{0, \dots, q-1\}$, the neighbors of $w$ are exactly $v_A = v_{i_A, 2j + 1 - r_A}$ and $ v_B = v_{i_B, 2j + 1 - r_B}$. We have $v_A \in I_2$ and $v_B \in J_2$. Since $V_{i_A, r_A} \subseteq I_1$ and $I_1$ is independent in $G$, we have $v_A \not\in I_1$. Similarly, $v_B \not\in J_1$. Since $I_2$ and $J_1$ are disjoint, we have $v_A \not\in J_1$, and similarly $v_B \not\in I_1$. Since $I_2$ and $J_1$ are not adjacent, we have $w \not\in J_1$, and similarly $w \not\in I_1$. However, this means that $I_1 \cup J_1$ is not a canonical solution since neither $w$ nor its two neighbors belong to $I_1 \cup J_1$. This is a contradiction, hence the claim is proven.
\end{proof}
\Cref{lem:is_thin_rects} follows immediately from the last claim.




\subsection{Lower bound on thin rectangle circuits}
In this subsection we prove \Cref{thm:main_is}. At this point we can abstract away from the exact structure of $G$. The proof will depend only on the thin rectangle property proven in the previous subsection. The proof will closely follow the proof of \cite[Lemma 2.18]{Jukna2023}.

Let $\Gamma$ be any tropical circuit that calculates $IS_G$. For any node $w$ of $\Gamma$, we define its \emph{below} $B_w$ as the polynomial calculated by $\Gamma$ if we designate $w$ to be its output node. Intuitively, $B_w$ captures the contribution of the subcircuit rooted at $w$ to the output of the whole calculation.

In a similar spirit, we would like to define the \emph{above} of a node $w$ as the contribution of the remainder of the circuit. Note that $B_w$ does not need to be contained in $IS_G$, however, assuming the output node is reachable from $w$, for at least one monomial $m$, we have $m \cdot B_w \subseteq IS_G$. We will define $A_w$ as the sum of all such monomials, that is
%$$
%    A_w = \sum \left\{ m \mid m \cdot B_w \subseteq IS_G \right\}.
%$$
$$
    A_w = \sum_{m:m \cdot B_w \subseteq IS_G} m.
$$
Naturally, $A_w \cdot B_w \subseteq IS_G$ for every node $w$ from which the output node is reachable.

Every monomial of each $A_w, B_w$ represents an independent set in $G$. Let $A_w^*$ denote the useful monomials of $A_w$, that is, monomials $m \in A_w^*$ such that there exist $m' \in B_w$ for which $m \cdot m'$ represents a canonical solution. We define $B_w^*$ analogously. As proven by \Cref{lem:is_thin_rects}, for every node $w$, one of $A_w^*, B_w^*$ must be of size at most $1$. If $|A_w^*| \leq 1$, we will say that $A_w \cdot B_w$ is \emph{$A$-thin}, otherwise, we will say that it is \emph{$B$-thin}.

If $w$ is an output node, then we have $A_w = \{1\}, B_w = IS_G$, so $A_w \cdot B_w$ is $A$-thin. If $w$ is an input node labeled $x$, we have $B_w = \{x\}$, therefore $A_w \cdot B_w$ is $B$-thin. Let $w_\ell, \dots, w_1$ be any path in $\Gamma$ where $w_1$ is an output node and $w_\ell$ is some input node. Based on our observations, there is some edge on this path $(w_{i+1}, w_i)$ such that $w_i$ is $A$-thin and $w_{i + 1}$ is $B$-thin, and so $|A_{w_i}^* \cdot B_{w_{i+1}}^*| \leq 1$.

Let $IS_G^*$ be the polynomial containing all monomials of $IS_G$ representing a canonical solution. The following claim shows that all canonical solutions belonging to the rectangles described above can be propagated along such input to output paths of $\Gamma$. We will use the term predecessor to refer to in-neighbors of a node of the circuit.

\begin{claim}
	Let $w$ be a node of $\Gamma$ and let $c \in IS_G^*$ be such that $c = a \cdot b$ for some $a \in A_w^*, b \in B_w^*$.
	Then, for at least one predecessor $u$ of $w$, we have $c \in A_u^* \cdot B_u^*$.
	Moreover, if $w$ is a + gate, then the above holds for both predecessors of $w$.
\end{claim}

\begin{proof}
	Let $u, v$ be the predecessors of $w$.
	First, consider the case where $w$ is a max gate, that is, $B_w = B_u + B_v$. Then, either $b \in B_u$ or $b \in B_v$. W.l.o.g. assume the former.
    We have $a \cdot B_u \subseteq a \cdot B_w \subseteq IS_G$, hence $a \in A_u$. Therefore, $b \in B_u^*$, $a \in A_u^*$, and hence $c \in A_u^* \cdot B_u^*$ as desired.
    
	Second, consider the case where $w$ is a + gate, that is, $B_w = B_u \cdot B_v$. Then, we have $b = b_v \cdot b_u$ for some $b_v \in B_v$ and $b_u \in B_u$.
    We have $a \cdot b_v \cdot B_u \subseteq a \cdot B_v \cdot B_u \subseteq IS_G$, hence $a \cdot b_v \in A_u$. Therefore, $b_u \in B_u^*$, $a \cdot b_v \in A_u^*$, and hence $c \in A_u^* \cdot B_u^*$. By a symmetric argument, $c \in A_v^* \cdot B_v^*$.
\end{proof}


Now, consider the following process. We fix a canonical solution $c \in IS_G^*$ and start at output node $w_1$. If we are currently in a max node $w_i$, we move towards the predecessor $w_{i+1}$ given by the claim. If we are currently in a + node, we move towards the predecessor $w_{i+1}$ for which $|B_{w_{i+1}}^*|$ is larger. Such process terminates at an input node and produces a path $w_\ell, \dots, w_1$. As argued before, for some $i \in [\ell - 1]$, we have $|A_{w_i}^* \cdot B_{w_{i+1}}^*| \leq 1$.

We have $c \in A_{w_{i+1}}^* \cdot B_{w_{i+1}}^*$. If $w_i$ is a max node, we have $B_{w_{i+1}} \subseteq B_{w_i}$, hence $A_{w_{i+1}}^* \subseteq A_{w_i}^*$, hence $c \in A_{w_i}^* \cdot B_{w_{i+1}}^*$, and hence $A_{w_i}^* \cdot B_{w_{i+1}}^* = \{c\}$.

If $w_i$ is a + node, then we additionally have $c \in A_v^* \cdot B_v^*$ where $v$ is the predecessor of $w_i$ other than $w_{i+1}$. Given the way we chose $w_{i+1}$, we have $|B_v^*| \leq |B_{w_{i+1}}^*| \leq 1$. Since $B_{w_i}^* \subseteq B_{w_{i+1}}^* \cdot B_v^*$, we have $|B_{w_i}^*| \leq 1$. Thus, $A_{w_i}^* \cdot B_{w_i}^* = \{c\}$.

Repeating this process for each canonical $c$ creates a mapping from $IS^*_G$ to the set $V(G) \cup E(G)$. If we map $c$ to a vertex $w$, then we have a guarantee that $A_w^* \cdot B_w^* = \{c\}$. If we map $c$ to an edge $(u,w)$, we have a guarantee that $A_w^* \cdot B_u^* = \{c\}$. This implies that the mapping is injective, and hence $2^k = |IS^*_G| \leq |V(G)| + |E(G)| \leq 3|V(G)|$. This finishes the proof of the lemma.


